\documentclass[twocolumn]{iopjournal}

\usepackage{amsmath,amssymb,amsfonts}

\fancyhfoffset[L]{0mm}
\setlength{\marginparwidth}{0pt}
\setlength{\marginparsep}{0pt}

% ==========================================
% ZENODO PREPRINT HEADERS (ACTIVE)
% ==========================================
\renewcommand{\articletype}[1]{%
  {\vspace*{-8mm}\noindent \Large \sf Preprint}\par
  \vspace*{6mm}\noindent{\scriptsize \sf{\bfseries \MakeUppercase{#1}}}}

\fancyhead[L]{}
\fancyhead[R]{}

\begin{document}

\articletype{Supplementary Material}

\title{The Energy Functional of the\\
Constitutive Law $\mu(x) = x/(x+1)$}

\author{Stephen Nelson}

\affil{Field Dynamics Research}

\ead{stephen@fielddynamics.org}

\begin{abstract}
In the AQUAL field theory of Bekenstein and Milgrom the constitutive
law $\mu$ enters through a kinetic function $F$, with
$F'(y) = \mu(\sqrt y)$ and $y = |\nabla\Phi|^2/a_0^2$. For
$\mu(x) = x/(x+1)$ we show that $F$ is uniquely
$F(y) = y - 2\sqrt y + 2\ln(1 + \sqrt y)$ once its irrelevant
additive constant is fixed, and give its convergent expansions at
both ends: Newtonian, $F = y - 2\sqrt y + \ln y + O(y^{-1/2})$, and
deep-MOND, $F = \tfrac23 y^{3/2} - \tfrac12 y^2 + O(y^{5/2})$,
where the leading term arises from a cancellation among all three
terms. As a function of the field gradient, the energy density is
twice continuously differentiable and strictly convex, with Hessian
eigenvalues $2\mu$ and $2(x\mu)'$. The field equation is therefore
elliptic wherever $\nabla\Phi \neq 0$, with the two eigenvalues
within a factor of two of each other, and on a bounded domain with
prescribed boundary values it has at most one solution.
\end{abstract}

% ----------------------------------------------------------------------

\section*{Setting}

The AQUAL Lagrangian~\cite{BekensteinMilgrom1984} is
\begin{equation}
  \mathcal{L} = -\frac{a_0^2}{8\pi G}\,
  F\!\left(\frac{|\nabla\Phi|^2}{a_0^2}\right) - \rho\,\Phi .
  \label{eq:L}
\end{equation}
Its Euler--Lagrange equation is
\begin{equation}
\begin{aligned}
  &\nabla\cdot\big[\mu(|\nabla\Phi|/a_0)\,\nabla\Phi\big]
  = 4\pi G\rho, \\
  &\mu(\sqrt y) = F'(y).
\end{aligned}
  \label{eq:field}
\end{equation}
An additive constant in $F$ leaves Eq.~\eqref{eq:field} unchanged,
and we fix it by $F(0) = 0$. Throughout,
\begin{equation}
  \mu(x) = \frac{x}{x+1}.
  \label{eq:mu}
\end{equation}

% ----------------------------------------------------------------------

\section*{The Kinetic Function}

\paragraph{Theorem 1.} \textit{The unique $F \in C^1[0,\infty)$
with $F'(y) = \mu(\sqrt y)$ and $F(0) = 0$ is}
\begin{equation}
  \boxed{F(y) = y - 2\sqrt y + 2\ln\!\big(1 + \sqrt y\big).}
  \label{eq:F}
\end{equation}

\paragraph{Proof.} Two such functions differ by a constant, which
$F(0) = 0$ removes. With $u = \sqrt t$,
\begin{equation}
\begin{aligned}
  F(y) &= \int_0^y \frac{\sqrt t}{1 + \sqrt t}\,dt \\
  &= 2\int_0^{\sqrt y}\Big(u - 1 + \frac{1}{1+u}\Big)\,du,
\end{aligned}
\end{equation}
since $u^2/(1+u) = u - 1 + 1/(1+u)$, and this integrates to
Eq.~\eqref{eq:F}. Conversely, with $u = \sqrt y$,
$dF/du = 2u - 2 + 2/(1+u) = 2u^2/(1+u)$, so
$F'(y) = (dF/du)/(2u) = u/(1+u)$, which extends continuously to
$y = 0$. \hfill $\square$

\paragraph{Corollary 1 (the two limits).} \textit{For
$0 \le y \le 1$,}
\begin{equation}
\begin{aligned}
  F(y) &= 2\sum_{k=3}^{\infty} \frac{(-1)^{k+1}}{k}\, y^{k/2} \\
       &= \tfrac23\, y^{3/2} - \tfrac12\, y^2 + \tfrac25\, y^{5/2}
         - \cdots,
\end{aligned}
  \label{eq:small}
\end{equation}
\textit{and for $y > 1$,}
\begin{equation}
  F(y) = y - 2\sqrt y + \ln y
       + 2\sum_{k=1}^{\infty} \frac{(-1)^{k+1}}{k}\, y^{-k/2}.
  \label{eq:large}
\end{equation}

\paragraph{Proof.} Insert
$\ln(1+u) = \sum_{k\ge1} (-1)^{k+1}u^k/k$ for $0 \le u \le 1$; the
terms $k = 1, 2$ cancel $-2u$ and $u^2$. For $u > 1$ write
$\ln(1+u) = \ln u + \ln(1 + 1/u)$ and expand the second logarithm.
\hfill $\square$

\medskip
\noindent
At high acceleration $F'(y) \to 1$ and Eq.~\eqref{eq:field} is
Poisson's equation. At low acceleration the leading term
$\tfrac23 y^{3/2}$ is the deep-MOND kinetic
function~\cite{BekensteinMilgrom1984}, for which
$\mu(x) \simeq x$ and the field equation becomes
$\nabla\cdot(|\nabla\Phi|\nabla\Phi) = 4\pi G a_0 \rho$. No single
term of Eq.~\eqref{eq:F} produces it: the orders $y^{1/2}$ and $y$
cancel among all three.

% ----------------------------------------------------------------------

\section*{Convexity and Ellipticity}

Write $q = \nabla\Phi/a_0 \in \mathbb{R}^3$ and
$W(q) = F(|q|^2)$, so that the energy density is
$(a_0^2/8\pi G)\,W(q) + \rho\Phi$ and Eq.~\eqref{eq:field} reads
$\nabla\cdot A(\nabla\Phi) = 4\pi G\rho$ with
$A(p) = \mu(|p|/a_0)\,p = \tfrac{a_0}{2}\nabla W(p/a_0)$.

\paragraph{Theorem 2.} \textit{$W \in C^2(\mathbb{R}^3)$, smooth
away from $q = 0$. Its Hessian at $q \neq 0$, with $x = |q|$ and
$\hat q = q/x$, is}
\begin{equation}
  \nabla^2 W = 2\mu(x)\,\big(I - \hat q\hat q^{\,T}\big)
             + 2\,(x\mu)'(x)\,\hat q\hat q^{\,T},
  \label{eq:hess}
\end{equation}
\textit{with $2\mu(x) = 2x/(1+x)$ twice and
$2(x\mu)'(x) = 2x(x+2)/(x+1)^2$ once, both positive, and
$\nabla^2 W(0) = 0$. Hence $W$ is strictly convex, $A$ is strictly
monotone, $(A(p) - A(p'))\cdot(p - p') > 0$ for $p \neq p'$, and
the field equation is elliptic wherever $\nabla\Phi \neq 0$, with}
\begin{equation}
  1 < \frac{(x\mu)'}{\mu} = \frac{x+2}{x+1} \le 2 .
  \label{eq:ratio}
\end{equation}

\paragraph{Proof.} $W(q) = g(|q|)$ with
$g(r) = F(r^2) = r^2 - 2r + 2\ln(1+r)$, so
$g'(r) = 2r^2/(1+r) = 2r\mu(r)$ and $g''(r) = 2(r\mu)'(r)
= 2r(r+2)/(1+r)^2$. For a radial function,
$\nabla W = g'(r)\,\hat q$ and
$\nabla^2 W = (g'/r)(I - \hat q\hat q^{\,T}) + g''\,\hat q\hat q^{\,T}$,
which is Eq.~\eqref{eq:hess}. As $r \to 0$, $g'/r = 2\mu(r) \to 0$
and $g'' \to 0$, so $\nabla W$ and $\nabla^2 W$ extend continuously
to $q = 0$ with value $0$, and $W \in C^2$; it is smooth for
$q \neq 0$. For $p \neq p'$ let
$\varphi(t) = W(p' + t(p - p'))$. Then
$\varphi''(t) = (p - p')^T\,\nabla^2 W\,(p - p') > 0$ except where
$p' + t(p - p') = 0$, which happens for at most one $t$, so $W$ is
strictly convex and
$\varphi'(1) - \varphi'(0) = \int_0^1 \varphi'' > 0$, which is the
monotonicity of $\nabla W$ and hence of $A$. The eigenvalues of
Eq.~\eqref{eq:hess} are those of the principal symbol of
$\nabla\cdot A(\nabla\Phi)$, up to the factor $2$, and
Eq.~\eqref{eq:ratio} follows from Eq.~\eqref{eq:mu}. \hfill
$\square$

\medskip
\noindent
At $\nabla\Phi = 0$ both eigenvalues vanish: the equation is
degenerate there, as the deep-MOND form
$\nabla\cdot(|\nabla\Phi|\nabla\Phi)$ already is.

\paragraph{Corollary 2 (uniqueness).} \textit{Let
$\Omega \subset \mathbb{R}^3$ be a bounded connected domain with
piecewise smooth boundary. Two solutions
$\Phi_1, \Phi_2 \in C^2(\Omega) \cap C^1(\bar\Omega)$ of
Eq.~\eqref{eq:field} with $\Phi_1 = \Phi_2$ on $\partial\Omega$
coincide.}

\paragraph{Proof.} $A$ is $C^1$, including at $p = 0$, where
$A(p) = |p|\,p/(a_0 + |p|)$ has derivative $0$. Let
$w = \Phi_1 - \Phi_2$, which vanishes on $\partial\Omega$. By the
divergence theorem and Eq.~\eqref{eq:field},
\begin{equation}
  \int_\Omega \big(A(\nabla\Phi_1) - A(\nabla\Phi_2)\big)
  \cdot\nabla w \, d^3x = 0 .
\end{equation}
The integrand is continuous and, by Theorem 2, positive wherever
$\nabla\Phi_1 \neq \nabla\Phi_2$. Hence $\nabla w = 0$ on $\Omega$,
$w$ is constant, and $w = 0$. \hfill $\square$

\medskip
\noindent
On all of space, Bekenstein and Milgrom showed that the condition
$(x\mu)' > 0$ makes the field equation elliptic and its solution
with $\nabla\Phi \to 0$ at infinity unique~\cite{BekensteinMilgrom1984};
by Eq.~\eqref{eq:ratio} this $\mu$ satisfies it.

% ----------------------------------------------------------------------

\section*{Summary}

The constitutive law $\mu(x) = x/(x+1)$ has the closed-form kinetic
function $F(y) = y - 2\sqrt y + 2\ln(1 + \sqrt y)$, unique up to the
constant that the field equation cannot see. It reduces to
Poisson's equation at high acceleration and to the deep-MOND
equation at low acceleration, through convergent expansions given in
full. As a function of the field gradient the energy is $C^2$ and
strictly convex, so the field equation is elliptic away from
$\nabla\Phi = 0$ and its solutions are unique for given boundary
values.

% ----------------------------------------------------------------------

\begin{thebibliography}{9}

\bibitem{BekensteinMilgrom1984}
J.~Bekenstein and M.~Milgrom,
\textit{Does the missing mass problem signal the breakdown of
Newtonian gravity?},
Astrophys.\ J.\ \textbf{286} (1984) 7.

\end{thebibliography}

\end{document}
